\section{Rangkuman}
\begin{itemize}
    \item GOST (Magma) adalah \textit{block cipher} standar Rusia dengan blok 64 bit, kunci 256 bit, jaringan Feistel 32 putaran, dan 8 kunci internal $K_1,\dots,K_8$ yang dijadwalkan berulang.
    \item Fungsi $f$ GOST adalah penjumlahan modulo $2^{32}$ dengan kunci putaran, substitusi 8 kotak-S ($4 \times 4$ bit), lalu pergeseran sirkuler kiri 11 bit.
    \item Dibanding DES, GOST memiliki panjang kunci yang jauh lebih besar (256 vs 56 bit), jumlah putaran yang lebih banyak (32 vs 16), dan kotak-S yang lebih sederhana ($4 \times 4$ vs $6 \times 4$). Keunggulan itu berlaku pada pencarian kunci menyeluruh dan pada efisiensi implementasi, \emph{bukan} pada ketahanan terhadap kriptanalisis struktural.
    \item Kriptanalisis atas GOST penuh 32 putaran telah menurunkan kompleksitas serangan dari $2^{256}$ menjadi sekitar $2^{178}$. Karena setiap serangan itu menuntut data paling sedikit $2^{32}$ blok --- tepat pada batas rekeying akibat paradoks ulang tahun --- GOST hanya dapat diserang bila kunci dipakai melampaui batas tersebut.
    \item Setiap \textit{block cipher} berukuran blok $n$ bit hanya boleh dipakai untuk mengenkripsi sekitar $2^{n/2}$ blok dengan satu kunci. Untuk blok 64 bit, batas itu adalah $2^{32}$ blok atau sekitar 32 GiB, dan berlaku bagi GOST maupun Blowfish.
    \item RC5 (Rivest, 1994) memiliki parameter fleksibel: ukuran kata $w \in \{16,32,64\}$, jumlah putaran $r \in \{0,\dots,255\}$, dan panjang kunci $b \in \{0,\dots,255\}$ byte.
    \item Jadwal kunci RC5 menyalin kunci ke larik $L$ secara \textit{little-endian}, menginisialisasi $S[0] = P_w$ dan $S[i] = S[i-1] + Q_w$, lalu mencampur $L$ dan $S$ sebanyak $n = 3 \cdot \max(t,c)$ langkah dengan $t = 2r+2$. Konstanta $P_w$ dan $Q_w$ diturunkan dari bilangan $e$ dan $\phi$ sehingga tidak dapat disembunyikan pintu belakang di dalamnya.
    \item Satu putaran RC5 adalah $A \leftarrow \mathrm{ROTL}(A \oplus B, B) + S_{2i}$ dan $B \leftarrow \mathrm{ROTL}(B \oplus A, A) + S_{2i+1}$, memakai $2r+2$ kunci putaran dengan operasi XOR, penjumlahan modulo $2^w$, dan rotasi yang besarnya ditentukan oleh data itu sendiri.
    \item RC6 memperluas RC5 ke blok 128 bit dengan kunci hingga 2040 bit dan menambahkan perkalian modulo $2^w$ untuk mempercepat difusi.
    \item Blowfish memakai larik $P$ berisi 18 kata dan empat kotak-S berisi 256 entri 32 bit yang seluruhnya dibangkitkan dari digit $\pi$ pada saat kunci dipasang. Twofish menggantinya dengan empat kotak-S $8 \times 8$ yang bergantung pada kunci, sebuah matriks MDS, dan transformasi PHT.
    \item Serpent adalah jaringan substitusi--permutasi murni dengan 32 putaran dan delapan kotak-S $4 \times 4$ bit; MARS adalah Feistel tak seimbang yang membagi 32 putarannya menjadi delapan putaran \textit{forward mixing}, enam belas putaran inti, dan delapan putaran \textit{backward mixing}.
\end{itemize}
